\section{Experimental design}\label{sec:design}

\subsection{Instances}
\paragraph{CVRPLIB.} We use 74 instances from CVRPLIB \citep{cvrplib}: set A (26 instances, 31--68 customers) and B (23, 30--77) with the classical definitions of Augerat et al.\ \citep{augerat1995}, set CMT (4, 50--120) of \citet{christofides1979}, and set X (21, 100--194) of \citet{uchoa2017}. The reference cost is the optimal or best-known value distributed with each instance.

\paragraph{Synthetic instances.}
To test generalization beyond CVRPLIB and to tune hyperparameters on instances disjoint from the evaluation set, we generate two synthetic sets with the X-instance generator of \citet{uchoa2017}, re-implemented as an in-process function and validated byte for byte against the official script (10 of 10 instances identical). \emph{Tuning set}: 24 instances (8 each with 50, 100 and 150 customers; seed 20261008). \emph{Validation set}: 30 instances (10 each with 100, 200 and 300 customers; seed 20261009). The combination of depot position, customer distribution, demand type and route length of each instance is sampled uniformly from the 378 combinations of the generator. The reference cost of each synthetic instance is the best of two runs of HGS-CVRP \citep{vidal2022} (seeds 1 and 2, 0.2~s per customer); it is not a proven optimum, and gaps are relative to it. As for set X, the number of routes of the reference solution is a lower bound.

\subsection{Budgets}\label{sec:budgets}
In Studies 1 and 2 the budget is the number of \emph{operator applications} $B\in\{800,4000\}$, counted by a wrapper around each neighborhood's best-improvement function; the shake does not count. The production configuration of the reproduced method (20 steps of 40 inner iterations) uses exactly 800 applications. With five similar operators this is a sound unit, but with operators whose cost differs by almost two orders of magnitude (Study~3) it would treat the most expensive operator as free. In Studies 3 and 4 the budget is therefore the number of \emph{distance evaluations} performed inside operators, $E_1=600\,n^2$ and $E_2=5E_1$ for $n$ customers, counted by a proxy of the instance that serves distances from a precomputed matrix. The count is deterministic given the seed, and the cost per evaluation is roughly constant (0.35--0.57~$\mu$s for $n$ between 53 and 195), so it tracks time without its noise on the heterogeneous cores of our machine. $E_1$ approximates the effort of 800 applications of the original pool; as a consistency check, the original pool with random order and local-optimum termination reaches 1.98\% at $E_2$ here and 1.94\% at $B=4000$ in Study~2. The check precedes each application; the application in progress completes, so a run overshoots its budget by at most the cost of one application (on average 0.03--0.6\%).

\subsection{Search structure}\label{sec:structure}
All arms share one generic loop (Algorithm~\ref{alg:loop}) parameterized by three structural factors---the operator \emph{selection policy}, the \emph{acceptance} rule and the inner-loop \emph{termination}---plus the pool. The initial solution is the Clarke--Wright savings solution repaired to satisfy the route-count constraint; the shake applies one random move from a randomly chosen original neighborhood and is identical across arms. Two code paths reproduce the earlier structures exactly (identical costs on all compared runs): fixed order/\emph{best}/LO reproduces the classical GVNS, and uniform/\emph{shaken}/F40 reproduces the original Algorithm~3 loop with uniform selection.

\begin{algorithm}[t]
\caption{Generic search loop. $B$ is the budget; policy $\pi$, acceptance $\in\{\textit{best},\textit{shaken}\}$, termination $\in\{\textrm{LO},\textrm{F40},\textrm{R}r\}$.}\label{alg:loop}
\begin{algorithmic}[1]
\State $x\gets$ initial solution; $x^*\gets x$
\While{budget not exhausted}
  \State $x'\gets \textsc{Shake}(x)$;\; $r\gets x'$;\; $F\gets\emptyset$;\; $m\gets 0$;\; $f\gets 0$
  \While{budget not exhausted \textbf{and} not stop}
    \State \textbf{stop if} LO: $|F|=k$;\; F40: $m=40$;\; R$r$: $f=r$
    \State choose operator $j\sim\pi$ (LO: among $\{1..k\}\setminus F$; else among all)
    \State apply best-improvement move of $j$ to $r$; update $\pi$ with the reward
    \If{improved} $r\gets$ new solution;\; $F\gets\emptyset$;\; $f\gets 0$ \Else\; $F\gets F\cup\{j\}$;\; $f\gets f+1$ \EndIf
    \State $m\gets m+1$
  \EndWhile
  \If{(\textit{best}: $c(r)<c(x^*)$) \textbf{or} (\textit{shaken}: $c(r)<c(x')$)} $x\gets r$;\; $x^*\gets r$ if $c(r)<c(x^*)$ \EndIf
\EndWhile
\State \Return $x^*$
\end{algorithmic}
\end{algorithm}

\subsection{Selection policies}\label{sec:policies}
\textbf{Uniform} (U) draws an operator uniformly; \textbf{fixed order} (FO) always takes the lowest-index allowed operator, which is the classical VND; \textbf{fixed order by cost} (FOC) cycles through a list ordered by increasing cost per application (fixed before the runs). \textbf{UCB1} \citep{auer2002} uses the reward normalized by the maximum observed, trying each operator once and then maximizing the normalized mean plus $C\sqrt{2\ln s/n_j}$ ($C=2$); \textbf{$\varepsilon$-greedy} uses the sample mean with $\varepsilon=0.1$; \textbf{probability matching} (PM) \citep{thierens2005} keeps an exponential moving average $Q_j\leftarrow Q_j+\gamma(r-Q_j)$ ($\gamma=0.1$) and selects with $p_j=p_{\min}+(1-kp_{\min})Q_j/\sum_lQ_l$ ($p_{\min}=0.05$); \textbf{HG-UCB} is described in Section~\ref{sec:hgucb}. Classical policies receive the relative improvement $(c'-c'')/c'$ on success and zero otherwise; HG-UCB receives its own reward. In Studies 3 and 4 the three learning policies use a \emph{cost-aware} reward: the success reward is multiplied by $\phi=\mathrm{clip}(\bar g/g,0.1,10)$, where $g$ is the distance evaluations of the application and $\bar g$ their running mean, so cheap successes are worth more; failures carry no cost information (a limitation we return to). The \textbf{static mixture} samples operators from a fixed distribution tuned offline (Study~4). Hyperparameter defaults come from the literature and from the original study (Appendix~\ref{app:hyper}); Study~4 tunes them.

\subsection{Studies}
Table~\ref{tab:studies} summarizes the four studies. Each was pre-specified (hypotheses, arms, tests and family size) in a dated research log before its data were produced (Appendix~\ref{app:prereg}).

\begin{table}[t]
\caption{Overview of the studies.}\label{tab:studies}
\centering\small
\begin{tabular}{p{1.4cm}p{3.6cm}p{2.8cm}p{3.6cm}p{2.0cm}}
\toprule
Study & Question & Structure and pool & Budget; instances; seeds & Confirmatory tests \\
\midrule
1 & Does the bandit beat chance? & Original loop (F40); 5 operators & $B\in\{800,4000\}$; 74 CVRPLIB; 10 & 12 (Holm) \\
2 & Where do differences come from? AOS inside VND? & Order $\times$ acceptance $\times$ termination; 5 operators & $B\in\{800,4000\}$; 74; 10 & 28 (Holm) \\
3 & Does a larger heterogeneous pool change it? & LO and R5; 10 operators & $E_1,E_2$; 74; 10 & 34 (Holm) \\
4 & Does tuning, a tuned static mixture or new instances change it? & R5; 10 operators & $E_1,E_2$; 74 and 30 synthetic; 10 and 5 & 14 and 6 (Holm) \\
\bottomrule
\end{tabular}
\end{table}

\subsection{Statistical protocol}
The unit of analysis is the instance: for each arm and instance the gap to the reference cost is averaged over the seeds (10 for CVRPLIB, 5 for the synthetic validation set), and arms are compared in pairs on the same instances. The primary test is the two-sided Wilcoxon signed-rank test \citep{wilcoxon1945}, because the paired differences are not normal (Shapiro--Wilk $p<0.05$ in 32 of 34 audited contrasts); omnibus comparisons use the Friedman test; multiplicity is controlled by the Holm procedure \citep{holm1979} over a family fixed before each study (Table~\ref{tab:studies}). Every comparison reports the mean paired difference in percentage points with a 95\% bootstrap confidence interval over instances (10\,000 resamples), wins, losses and ties. Because ``no significant difference'' is not evidence of equivalence, claims of no difference are supported by equivalence tests \citep{lakens2017}: two arms are equivalent within $\delta$ when the 90\% bootstrap interval of the mean difference lies within $(-\delta,+\delta)$ (a two one-sided tests procedure at 5\%), with $\delta=0.10$~pp as the primary margin (smaller than every confirmed structural effect, 0.13--0.8~pp) and 0.05, 0.15 and 0.20~pp as sensitivity. As a check on pseudoreplication, we also report 95\% intervals from a two-level bootstrap (instances, then seeds within instance, keeping the pairing by seed), and the minimum detectable effect at 80\% power, $2.8\,\mathrm{SD}_d/\sqrt{n}$ with $\mathrm{SD}_d$ the standard deviation of the paired instance differences. These robustness analyses were chosen after the confirmatory tests and are labeled as such.

\subsection{Verification}
Four controls precede the studies. (i) \emph{Exact regression} between code versions: each new version of the search loop must reproduce, cost for cost, the runs of the previous version under identical seeds (48 of 48 comparisons for the loop used in Study~2 after rewriting it; 8 of 8 and 6 of 6 for the cost-counting and instrumented paths). (ii) \emph{Operator validation}: the five additional operators were validated against an independent brute-force enumeration of their neighborhoods (1\,600 tests including infeasible-move and route-count cases), and a mutation test confirmed that the test can fail (a deliberately degraded SWAP* was detected in 29 of 40 cases). (iii) \emph{Accounting}: operator applications or distance evaluations match the budget in every run, and every returned solution is feasible (no infeasible run in any study). (iv) \emph{External reference}: HGS-CVRP solves all 74 instances in all runs with a feasible solution. \pend{independent recomputation of the headline gaps from raw costs by a script that shares no code with the analysis scripts.}

\subsection{Computing environment}
All runs used one laptop (Intel Core i5-12450H, 12 threads, 7.8~GB visible to the Linux subsystem), Windows 11 with an Ubuntu 26.04 LTS WSL2 guest, Python 3.14.4 (NumPy 2.5.3, SciPy 1.18.1, pandas 3.0.5, river 0.26.1 for ADWIN), and up to eleven parallel worker processes. HGS-CVRP (commit \texttt{1a92795}) was compiled with g++ 15.2.0 (\texttt{-O3}). Because budgets are in operator applications or distance evaluations, results do not depend on this hardware; wall-clock times are reported only descriptively.
